\section{Simultaneous hook and marked-site trace bounds}
\label{ten:simultaneous}
This proof keeps two assertions simultaneously: a dimension bound for every representation, and a stronger bound for every permitted hook subdiagram with a sufficiently large marked complement. The positive-part penalty in the trace recursion controls the sum of child log dimensions.
\smallskip
All notation introduced below is local to this section. Traces are unnormalized, logarithms are natural, and a sweep on $2^d$ positions takes $d$ physical shuffles. The conventions are those of Section~\ref{sec:prelim}.

\subsection{Moment normalization and simultaneous targets}
We work over \(\mathbb C\), using unitary representations of symmetric groups. For a dyadic size $m=2^a$, use the sweep $B_m$ of Section~\ref{sec:prelim}. Partition indices of irreducibles are denoted \(\lambda\), the representations by \([\lambda]\), their dimensions by \(D_\lambda\), and traces in them by \(\operatorname{tr}_\lambda\). Use also the empty partition in size zero, of dimension 1. From this point through the end of this section, \(p\ge1\) denotes an integer square-moment index; Proposition~\ref{ten:simultaneous:main} later chooses \(p\) to be an absolute power of two. For \(\lambda\vdash m=2^a\), set
\[
 \begin{aligned}
 B_\lambda&=H_\lambda=\log D_\lambda,\\
 \mathcal L_\lambda&=\operatorname{tr}_\lambda Q_m^p
     =\operatorname{tr}_\lambda|B_m|^{2p},\\
 e_\lambda&=B_\lambda+\log\mathcal L_\lambda
     =E_m(\lambda;2p).
 \end{aligned}
\]
Use \(e_\lambda=-\infty\) if the trace is zero. \(B_\lambda\) will also be used in sizes other than powers of 2.

The general marked recurrence will be used at real Schatten order $2p$.
The simultaneous induction below needs its sparse estimate, a dimension
lower bound, and a comparison when the permitted hook shrinks. After
completing that induction, we give an alternative derivation of its
recurrence using weight projections and a trace inequality for powers
of two. That derivation reuses the ordinary representation-theoretic
conventions but supplies its own density and marked-trace argument.

The aim is to bound \(\mathcal L_\lambda\) for one fixed large \(p\). Near a single long row we do this directly using functions of a sparse set of labels. Otherwise we split a sweep in half, into independent operations on rows of a grid and then on the transversal columns, and recurse on traces. A graded tensor space handles diagrams in a small hook: it gives a weight-entropy bound from log dimension, and product states on rows and columns permit an overlap bound using the transversality. Distinct marked slots handle cells outside a hook. To keep the errors small relative to dimension the induction will track both log dimension on the hook and a budget for the marks. Constants in estimates below may increase from one bound to the next and are absolute once the indicated numerical parameters are fixed.

Take \(p\) to be a power of 2; its size will be fixed after all thresholds. Use \(\gamma=1/40,\ \eta=9/10\), and \(K_n=\lfloor n^\gamma\rfloor\). Define
\[
 J(t,n)=t\log(n/t),\qquad W(t,n)=t\log n
\]
(zero at \(t=0\)). Let \(n=2^d\) be sufficiently large, and split the sweep into first \(\lfloor d/2\rfloor\) and last \(\lceil d/2\rceil\) coordinates, with \(r,s\) the respective powers of 2. The first part \(X\) acts by copies of \(B_r\) within the \(s\) rows of the grid, the second \(Y\) by copies of \(B_s\) in the \(r\) columns: averages factor across the disjoint groups in each part. So \(B_n=YX\). Index these groups together by \(g\), of sizes \(m_g\).

Suppose \(\lambda\vdash n\) has \(\mathcal L_\lambda>0\), and choose \emph{any} subpartition \(\mu\subseteq\lambda\) of size \(M=n-t\) fitting the \(K_n\)-hook. Proposition~\ref{s:general-marked-recurrence}, with real Schatten order $2p$
and hook parameter $q=K_n$, gives
\begin{equation}
 e_\lambda+J(t,n)
 \ \le\ \max\ \Big\{
 \sum_g \big(e_{\alpha_g}+J(h_g,m_g)\big)
 -\big(\sum_g B_{\delta_g}-B_\mu\big)_+ \Big\}+ C(t+n^\eta).           \label{ten:simultaneous:eq:6}
\end{equation}
Here one can take the max over choices with integers \(0\le h_g\le m_g\) summing to \(t\) in each of the two branches (rows, columns);
\(\alpha_g\vdash m_g,\ \mathcal L_{\alpha_g}>0\); and \(\delta_g\subseteq\alpha_g\) of size \(m_g-h_g\) fitting the \(K_n\)-hook. The constants and size threshold do not depend on \(p\).


Indeed its error is $O(t+n^{9/10})$, since
$(r+s)(K_n^2+n^{1/4})\log(n+1)=O(n^{9/10})$.
The alternative weight-space proof appears at the end of this section.

\subsection{The simultaneous estimates}


Fix \(u=10^{-5}\), \(\delta=u/4\); the sparse and dimension estimates below will use this value of \(\delta\). Call \(t\) high in size \(n\) if \(t>0\) and \(\log(n/t)\le u\log n\). Use bit-length levels \(d\le 2^l D\), \(l\ge0\), with an absolute integer \(D\) sufficiently large. Set
\[
 A_0=1/4,\qquad b_0=\gamma/100,\qquad
 A_l-A_{l-1}=b_l-b_{l-1}=(2^{l-1}D)^{-1/2}\quad(l\ge1).
\]
Take \(D\) large enough that the sum of increments is less than \(b_0\); in particular \(A_l<1/2,\ 3u<b_l\le 2b_0<A_0\gamma/2\). Further sufficiently-large requirements below on \(D\) will be independent of the eventual moment exponent.

\begin{proposition}[Simultaneous dimension and marked-hook estimates]
\label{ten:simultaneous:main}
There are an absolute threshold \(D\) and, after \(D\) is fixed, an
absolute power of two \(p\ge1\) such that the following bounds hold
at every level \(l\ge0\) and every size \(n=2^d\) with \(1\le d\le2^lD\).
For every \(\lambda\vdash n\) with positive trace, (G) holds. For every
such \(\lambda\) and every \(\mu\subseteq\lambda\) with \(\mu_{K_n+1}\le K_n\),
put \(t=n-|\mu|\). If \(t>0\) and \(\log(n/t)\le u\log n\), then (H) holds:
\begin{equation}
 \begin{array}{ll}
 \text{(G)} & e_\lambda\le A_l B_\lambda,\\
 \text{(H)} & e_\lambda+J(t,n)\le A_l B_\mu+b_l W(t,n).
 \end{array}                                                         \label{ten:simultaneous:eq:8}
\end{equation}
\end{proposition}

\subsection{Sparse harmonic-tuple coupling}
Here first allow any fixed $0<\delta<1/2$; the simultaneous induction
uses the value chosen above.

There is a \(p_0=p_0(\delta)\) such that, for any integer \(p\ge p_0\) and sufficiently large \(m=2^a\) (threshold depending only on \(\delta\)),
\begin{equation}
 1\le k=m-\lambda_1\le m^{1-\delta}
 \quad\Longrightarrow\quad e_\lambda\le -2k\log m.                     \label{ten:simultaneous:eq:1}
\end{equation}
For this consider the space of functions on ordered injections \(x=(x_1,\ldots,x_k)\) of positions, with the uniform inner product, under permutation of positions (\((\pi f)(x)=f(\pi^{-1}x)\)). Call a function harmonic here if summing over the choices of any one coordinate, with the others fixed, gives zero. These functions form an invariant subspace: it is the orthogonal complement of the span of the functions independent of at least one coordinate. Applying the branching rule \cite[Theorem~5.8]{VershikOkounkov2005} and Frobenius reciprocity \cite[Theorem~4.33]{etingof}, \([\lambda]\) with first row \(m-k\) occurs on the injection space (the stabilizer is \(S_{m-k}\)), but does not occur on \((k-1)\)-injections, since it does not contain the subpartition \((m-k+1)\). Its isotypic subspace is therefore harmonic.

Let \(T=(B_m^* B_m)^p\) act on these functions. Evaluation of \(g=T f\) is the expectation of \(f\) by evolving positions through \(p\) blocks, each consisting of two sweeps (one in reverse order; reversing the factors for action on positions still has the two-sweep property). We keep layers as separate random layers even where a projection is repeated. For disjoint injections \(x,y\) (all \(2k\) positions distinct, possible for these sizes when \(m\) is sufficiently large), consider the alternating sum \(\Delta g(x,y)\) of \(g\) on the \(2^k\) injections choosing \(x_i\) or \(y_i\) for each \(i\), with sign \((-1)^{\#\{i:y_i\text{ chosen}\}}\).

Couple all these evolutions by a skeleton with two alternative paths per label \(i\), initially at \(x_i,y_i\). They can coalesce within a label, in which case we use one path thereafter. Paths of distinct labels stay distinct. At a layer, call a label clean if none of its paths occupies a switch pair shared with a path of a different label. For each clean label choose a fresh uniform \emph{destination bit} in the active coordinate, the same for both its paths, directing each to that endpoint of its pair. (This can direct both to the same endpoint.) For pairs involving unclean labels use the usual swap coins. All coins except the sharing just specified are fresh and independent. Pairs shared across labels always use a common usual swap, so preserve distinctness. For any fixed selection of alternatives, conditionally on the skeleton so far its paths have exactly the correct one-step transition: sharing a destination bit within a clean label causes no required correlation to fail, since only one of its paths is selected; among the pairs seen by the selection there are independent coins and a shared pair makes complementary outputs. Consequently from any fixed skeleton at a block start the trajectory of a fixed selection has the standard shuffle-path law.

Once any label coalesces, the alternating sum of the final \(f\) evaluations cancels pointwise in this coupling (its expectation is \(\Delta g\)), by toggling that label's choice. To have no coalescence through a block starting with none, every label must be unclean at least once in the \emph{last} sweep of the block: clean updates all through that sweep give identical bits in every coordinate for its paths. In such a skeleton, a uniformly random selection of alternatives has in expectation at least \(k/4\) distinct labels that participate in selected-path pair collisions in this sweep. Indeed for each label take a witnessing shared pair; selection of the two paths at that pair has probability \(1/4\) since they belong to distinct labels. In particular some selection has at least \(k/4\) participants. We union bound this last event over the \(2^k\) fixed selections, without conditioning their laws on what happens during the block.

For one fixed selection, we can do the trajectory probability calculation in its standard path law starting from deterministic positions conditionally at the block start. At the start of the last sweep the occupied subset satisfies
\[
 \mathbb P(A\text{ all occupied})\le (k/m)^{|A|}
\]
for any set \(A\). To see this, the same product bound using individual marginal occupation probabilities holds initially from deterministic positions and is preserved by a layer. Pull \(A\) back through the independent swaps; the expected product of old marginals factors by pairs. A single endpoint in \(A\) yields the averaged marginal, and both endpoints yield the old product at most the square of the average. After the first sweep the averaged marginals are all \(k/m\), by averaging along every coordinate. Independently of this subset, full shuffle coins in the last sweep define a graph on the \(m\) positions at the start of that sweep. Trace paths starting at all positions, connecting starting positions whose paths occupy a switch pair at any layer just before its update. The graph has maximum degree at most \(a\). Participants are exactly the occupied vertices with an occupied neighbor.

In a fixed such graph, the nontrivial connected components of the induced occupied graph, if they have \(v\ge k/4\) vertices total, give a collection of \(c\le v/2\) nontrivial connected sets. The number of possibilities for given \(v,c\) is at most \(\binom mc 2^v a^{2v}\): choose roots in order (say the least elements, sorted), sizes, and describe each connected set by a walk traversing a spanning tree from its root using \(2(\text{size}-1)\) steps. Apply the occupied-set bound to each such collection (with \(v\) distinct vertices). Since \((em/c)^c\) is increasing up to \(c=m\), we can bound the root count by \((2em/v)^{v/2}\). Summing in \(c\) (at most \(v\le 2^v\) terms), the probability for a fixed selection is at most
\[
 \sum_{v=\lceil k/4\rceil}^{k} [C a^2\sqrt{k/m}]^v
\]
for an absolute \(C\), using \(v\ge k/4\). For sufficiently large \(m\) and \(q=k/m\le m^{-\delta}\), the bracket is at most \(q^{1/4}\), giving sum at most \(2q^{k/16}\). Including the \(2^k\) selections, the conditional probability of no coalescence through the block is at most \(q^{k/32}\) (absorbing \(2^{k+1}\) for sufficiently large \(m\)). This bound repeats conditionally at each block. It gives
\[
 |\Delta g(x,y)|\le 2^k q^{pk/32}\|f\|_\infty .
\]

If \(f\) is harmonic so is \(g\). Average the alternating difference over a uniform \(y\) disjoint from \(x\). For a nonempty subset \(S\) of coordinates replaced, of size \(j\), summing \(g\) over distinct choices for these outside all of \(x\) can be done by inclusion-exclusion starting with choices just avoiding the positions of \(x\) in coordinates outside \(S\) (and still distinct). Expand \(\prod_{i\in S}(1-\mathbf 1_{\{y_i\in x_S\}})\) where \(x_S\) here is the set of corresponding positions. Terms with any unrestricted variable coordinate vanish on summing in that coordinate by harmonicity. The result is \((-1)^j\) times the sum with a permutation of the \(j\) positions \(x_S\) used, divided by \((m-k)_j\) for averaging. Thus the average alternating difference differs from \(g(x)\) by at most
\(\|g\|_\infty\sum_{j=1}^k (k)_j/(m-k)_j\), where \((m)_j=m(m-1)\cdots(m-j+1)\), with \((m)_0=1\). The sum is at most \(1/2\) for sufficiently large \(m\) here (bound by a geometric sum with ratio \(k/(m-2k)\)). We conclude that \(\|T f\|_\infty\le 2^{k+1} q^{pk/32}\|f\|_\infty\) on the harmonic subspace. Applied to eigenvectors, this bounds the eigenvalues on \([\lambda]\); \(T\) is positive semidefinite. Also \(D_\lambda\le m^k\) by its occurrence, so \(e_\lambda\le 2k\log m+(k+1)\log 2-(p\delta/32)k\log m\). We get \eqref{ten:simultaneous:eq:1}, for example taking \(p_0\ge 192/\delta\).

\subsection{Dimensions in the nonsparse range}
If \(\mathcal L_\lambda>0\), the length of \(\lambda\) is at most \(m/2\). In fact invariants for the first coordinate pair-swap subgroup \(S_2^{m/2}\) must occur, and induction of its trivial module builds the partition by \(m/2\) horizontal 2-strips. We claim for sufficiently large \(m\),
\begin{equation}
 \mathcal L_\lambda>0,\quad m-\lambda_1>m^{1-\delta}
 \quad\Longrightarrow\quad B_\lambda\ge c m^{1-\delta}                \label{ten:simultaneous:eq:2}
\end{equation}
with \(c>0\). If \(\lambda_1\ge m/2\), interleave the first row with one fixed standard order of the \(k=m-\lambda_1\) tail cells, ensuring row count at least tail count in every prefix; this suffices for all column conditions (a tail cell in column \(j\) appears no earlier than the \(j\)-th tail step). The ballot count is at least \(\binom mk/(m+1)\) (subtract \(\binom m{k-1}\) by reflection at the first bad prefix), and \(\binom mk\ge(m/k)^k\ge2^k\). Otherwise both maximum row and column are at most \(m/2\). If their maximum \(b\) is at least \(m/(4e)\) (here \(e\) is the base of natural logarithms), orient the partition, by transposing if needed, with row size \(b\) and take a subpartition with that row and tail size \(b\). The same interleavings have log count \(\Omega(m)\). Dimension does not decrease in extending a subpartition (extend tableaux), nor change in transposing. If the maximum is smaller, all hooks are at most \(m/(2e)\) and the hook formula gives dimension at least \(2^m\). This proves \eqref{ten:simultaneous:eq:2}.

\subsection{Shrinking the allowed hook}
We need a quantitative dimension estimate when shrinking the permitted hook. Let \(\theta\) have size \(M\le m\), let \(K=K_m\), and remove its cells with both indices \(>K\), obtaining \(\widehat\theta\) by removing \(v\) cells. Always \(B_\theta\ge B_{\widehat\theta}\). For sufficiently large \(m\),
\begin{equation}
 M\ge K^3\quad\Longrightarrow\quad
 B_\theta-B_{\widehat\theta}\ge (\gamma/2)v\log m.                     \label{ten:simultaneous:eq:7}
\end{equation}
For \(v>0\), write \(a,b\) for width and height of the removed region, taking \(a\ge b\) by symmetry. The dimension ratio by hooks has numerator \((M)_v\); the denominator consists of the removed hooks (at most \(2a\) each) and the inflation factors of the preserved hooks.

For the top \(K\) rows only columns \(K+j,\ 1\le j\le a\), have inflation. Let \(b_j\) be the decreasing (nonincreasing) heights added there. The hooks before adding are at least \(a-j+u,\ u=1,\ldots,K\), since the full rectangle there in the top rows was present. Per column the log inflations thus sum to at most \(\log\binom{K+b_j}{b_j}\le b_j(1+\log(1+K/b_j))\), just using the denominator bound \(u\). Summing by concavity gives at most \(v(1+\log(1+Ka/v))\). Another bound using \(\log(1+x)\le x\) is
\[
 \sum_j b_j\sum_{u=1}^K\frac1{a-j+u}
 \ \le\ \frac va\sum_{j,u}\frac1{a-j+u}
 \ \le\ v(2+\log(1+K/a)).
\]
The first bound here averages oppositely ordered sequences; the double sum is at most \(\sum_{i=1}^a 1/i+\log\binom{a+K-1}{a}\) by summing in \(u\) with an integral.

Similarly for the left \(K\) columns with \(b\) in place of \(a\): choosing the better of the concavity and opposite-order bounds gives log inflation per removed cell at most \(2+\log(1+K/\max(v/b,b))\). For the top rows we can just use \(2+\log(1+K/a)\). The total log denominator per removed cell is therefore at most
\[
 C+\log\!\big[ a \max(1,K/a)\max(1,K/\sqrt a)\big],
\]
using \(\max(b,v/b)\ge\sqrt v\ge\sqrt a\). The bracket is at most \(\max(M/K,K^2)\): \(a\le M/K\), and consider \(a\ge K^2,\ K\le a\le K^2,\ a\le K\) respectively. For \(M\ge K^3\), comparison with \((M/e)^v\) for the numerator proves a gain per cell of \(\log K-C-1\), giving \eqref{ten:simultaneous:eq:7}. Nondecrease without the size assumption is by tableau extension.

\subsection{Proof of the simultaneous estimates}

\begin{proof}[Proof of Proposition~\ref{ten:simultaneous:main}]
For level 0 we can, after choosing \(D\), take a power of 2 \(p\ge p_0\) sufficiently large and fixed. In each nontrivial irreducible the sweep norm is strictly less than 1: equality for a unit vector would require equality at every successive orthogonal projection, leaving the vector invariant under all hypercube edge transpositions, which generate the whole symmetric group by connectivity. Thus in the finitely many base sizes we can have \(e_\lambda\le -n\log n\) for all nontrivial types. In the trivial type \(e_\lambda=0\). This proves the base (\(J\le n\log n\), and use \(J\le uW\) in the high case for the trivial type).

Suppose \(2^{l-1}D<d\le 2^l D\), so the children are at most level \(l-1\) and their bit lengths are at least \(d/3\). Write \(A=A_{l-1}, b=b_{l-1}\). First prove (H). Apply \eqref{ten:simultaneous:eq:6} with the chosen \(\mu\). Consider any child choice there, writing \(m=m_g,\ h=h_g,\ \alpha=\alpha_g,\ \theta=\delta_g,\ S_g=B_\theta,\ W_g=h\log m\). We have \(\sum_g W_g=W(t,n)\), since each branch sums the counts to \(t\) and \(\log r+\log s=\log n\).

\begin{itemize}
\item If inherited \(h\) is high in size \(m\), truncate \(\theta\) to \(\widehat\theta\) as in \eqref{ten:simultaneous:eq:7}, adding \(v\) removed cells to the count \(h\). Use child (H) with \(\widehat\theta,h+v\), still high. Since \(J(h,m)-J(h+v,m)\le v\) (differentiate in the count), we get
\[
 e_\alpha+J(h,m)
 \le A S_g+b W_g+(b\log m+1)v-A(B_\theta-B_{\widehat\theta}).
\]
When \(m-h\ge K_m^3\) the extra is nonpositive by \eqref{ten:simultaneous:eq:7} and the strict bound on \(2b_0\), for sufficiently large \(D\). Otherwise \(v\le m^{3\gamma}\), \(h\ge m/2\), and dimension is nondecreasing, so the extra is at most \(2(b+1/\log m)m^{3\gamma-1}W_g\).
\end{itemize}

\begin{itemize}
\item Otherwise \(h<m^{1-u}\) (including \(0\)). If the child is sparse (\(m-\alpha_1\le m^{1-\delta}\)), then \(e_\alpha\le0\) by \eqref{ten:simultaneous:eq:1} or triviality. If not, (G) at child size gives \(e_\alpha\le A(S_g+W_g)\), since \(D_\alpha\le m^h D_\theta\) by the same tableau-region bound (\(f^{\alpha/\theta}\le h!\)). By \eqref{ten:simultaneous:eq:2}, \(B_\alpha\ge c m^{1-\delta}\), so \(W_g/S_g=O(m^{-(u-\delta)}\log m)\) for sufficiently large sizes by \(S_g\ge B_\alpha-W_g\); we absorb \(A W_g\) by this small increase in the coefficient on \(S_g\). In both cases we can bound \(J\) by
\[
 J(h,m)\le 3u W_g + h\log^+(\bar h/h),\qquad \bar h=tm/n,
\]
with zero last term at \(h=0\), \(\log^+ x=\max(0,\log x)\). Indeed split \(\log(m/h)=\log(n/t)+\log(\bar h/h)\) and \(\log n\le3\log m\). The last terms sum to at most \(2t/e\) over all children, since each is at most \(\bar h/e\) and the \(\bar h\)'s sum to \(t\) in each branch (this bound also allows including the other children).
\end{itemize}

Together the child sum in \eqref{ten:simultaneous:eq:6} is at most
\[
 (A+\epsilon_d)\sum_g S_g + (b+\epsilon_d)W(t,n)+2t/e
\]
where \(\epsilon_d\ge0\) is \(o(d^{-1/2})\) uniformly: the errors just used besides the dilution terms are bounded in the coefficients by \(C m^{3\gamma-1}\) and \(C m^{-(u-\delta)}\log m\), \(2^{d/3}\le m\le 2^d\). Since \(A+\epsilon_d\le1\), the positive-part penalty in \eqref{ten:simultaneous:eq:6} reduces the structural term to at most \((A+\epsilon_d)B_\mu\) (use \(aS-(S-B_\mu)_+\le a B_\mu\) for \(0\le a\le1\)). Also \(W(t,n)\ge n^{1-u}\log n\) here, so \((C(t+n^\eta)+2t/e)/W(t,n)=o(d^{-1/2})\). These coefficient errors are bounded together by the level increment for sufficiently large \(D\), proving (H).

To prove (G), the sparse or trivial case needs no recursion. Otherwise use \(\mu=\widehat\lambda\), the canonical hook truncation. If its removed count is high, apply (H) just proved, using \eqref{ten:simultaneous:eq:7} (its parameters \(M=m=n,\theta=\lambda\)) to absorb \(b_l W\) into the dimension gain times \(A_l\). If it is not high, use \eqref{ten:simultaneous:eq:6} once with only child (G), bounding each \(e_{\alpha_g}+J(h_g,m_g)\le A S_g+(A+1)W_g\). The penalty again controls the structural sum. Thus
\[
 e_\lambda\le A B_\lambda+(A+1)W(t,n)+ C(t+n^\eta)
\]
by dimension nondecrease. Here \(t<n^{1-u}\) and \eqref{ten:simultaneous:eq:2} applies to \(\lambda\), making the extra divided by \(B_\lambda\) at most \(C(n^{-(u-\delta)}\log n+n^{\eta-1+\delta})=o(d^{-1/2})\). This proves (G).

In these steps \(D\) can be chosen once: the sufficiently-large size thresholds hold also for \(m\ge2^{d/3}\) as needed, and the displayed little-oh errors (with absolute bounds independent of levels and \(p\ge p_0\)) can be made smaller, including the indicated sums, than \(d^{-1/2}\le(2^{l-1}D)^{-1/2}\) for all \(d>D\). Sizes already treated retain \eqref{ten:simultaneous:eq:8} when increasing the constants. This completes the simultaneous induction with an absolute fixed \(p\).

\end{proof}

\subsection{Full-deck amplification}


For a positive integer \(z\), Schatten H\"older gives \(\|B_n^z\|_{\rm HS}^2\le \operatorname{tr}((B_n^*B_n)^z)\) in each irreducible (use Schatten exponent \(2z\) on each factor). Take \(z=4p\). In the nonsparse positive-trace types, this trace is at most \(\mathcal L_\lambda^4\), so its product with \(D_\lambda\) is at most \(\exp(-B_\lambda)\) by (G) and \(A_l<1/2\). Their sum tends to zero by \eqref{ten:simultaneous:eq:2} and the Durfee-square bound on the number of partitions. In the nontrivial sparse types use just the trace at \(p\) as upper bound (sweep norm at most 1), giving products at most \(n^{-2k}\) by \eqref{ten:simultaneous:eq:1}, summable to \(o(1)\) since there are at most \(2^k\) tail partitions of size \(k\). Zero traces contribute nothing.

Shuffle distributions act the same on positions from any starting deck; uniform orders correspond to uniform permutations. At \(zd\) shuffles the position law is \(B_n^z\), \(n=2^d\). Finite group Plancherel (or the regular trace applied to the squared norm of the group-algebra difference from uniform) identifies \(n!\) times the coefficient-wise squared distance with \(\sum_{\lambda\ne(n)}D_\lambda\|B_n^z\|_{{\rm HS},\lambda}^2\). By Cauchy--Schwarz its total variation distance from uniform is at most one half the square root of this sum, which tends to zero. This gives the upper mixing bound.

\subsection{An alternative recurrence proof through weight projections}

For completeness, we now derive \eqref{ten:simultaneous:eq:6} by a
second overlap mechanism. The preceding induction used the general
full-isotypic recurrence. The argument here instead resolves the global
hook type into torus weights and extracts its entropy from those weights.
We retain the notation $B_\theta=\log D_\theta$, the grid and marked
counts in \eqref{ten:simultaneous:eq:6}, and the fixed power-of-two
square-moment index $p$.

\subsubsection{Graded tensor density estimates}


For an integer \(K\ge 1\) let \(V=E\oplus O\) be an orthogonal sum of an even and an odd space, each of dimension \(K\). On \(H_k=V^{\otimes k}\) use the signed slot permutations: when permuting homogeneous factors the sign is that of the reordering of the odd factors. This gives a unitary \(S_k\)-action (successive reorderings multiply the signs), commuting with \(U(E)\times U(O)\) acting diagonally over the slots. Reordering the slots by a signed permutation conjugates tensor products of parity-preserving slot operators to their ordinary reordered products: the sign correction commutes with operators preserving each slot's parity. Subgroup actions on disjoint subsets of slots become ordinary tensor products of the respective signed actions after grouping the subsets by such a reordering.

Let \(\Pi_\theta\) denote the \(S_k\)-isotypic projection and \(w_\theta\) the multiplicity in \(H_k\). Call the \(K\)-hook condition \(\theta_{K+1}\le K\). In the next estimates take \(n\ge 2,\ K,k\le n\). For a count vector \(u\) of sum \(k\), put
\[
 H(u)= k\log k-\sum_j u_j\log u_j
\]
with zero contributions at zero (also \(H=0\) if \(k=0\)). A weight subspace refers to the span of tensors with a given count vector in an orthonormal parity basis. With absolute constants \(C\) (which can be increased), the following hold:
\begin{itemize}
\item \(w_\theta>0\) exactly under the \(K\)-hook condition, and then \(w_\theta\le(n+1)^{3K^2}\).
\item In any parity basis the range of \(\Pi_\theta\) lies in the sum of weight subspaces with \(H(u)\ge B_\theta-C K^2\log(2n)\).
\item For each occurring \(\theta\) there is an average over parity-preserving density matrices \(\rho\) (positive semidefinite, trace 1) on \(V\) such that in the positive semidefinite order,
\begin{equation}
 \Pi_\theta\ \le\ \exp(B_\theta+C K^2\log(2n))\,\mathbb E\,\rho^{\otimes k}.       \label{ten:simultaneous:eq:3}
\end{equation}
\end{itemize}

For the first fact, decompose into subspaces with fixed basis counts. Each is an induced monomial representation from the word stabilizer with trivial characters on even label groups and sign characters on odd ones. By Pieri, summing the multiplicity over counts amounts to counting chains ending at \(\theta\) that add first \(K\) horizontal strips then \(K\) vertical strips (sizes can be zero). The even stages have at most \(K\) rows; the odd stages have cells below the top \(K\) rows only in the first \(K\) columns. This proves the necessity, and one can build a fitting \(\theta\) by its first \(K\) rows then by the columns below. The shapes (hence also strip sizes) are determined by at most \(K\) row lengths at each even stage and the first \(K\) row and first \(K\) column lengths at each odd stage, giving the bound.

Alternatively, at fixed numbers of even and odd slots the representation is induced, by varying their locations, from tensor powers of \(E\) and \(O\) in the individual fixed parity pattern (one can use all even first); the latter power is sign-twisted for the slot permutations. Schur--Weyl in that pattern then gives \(U(E)\times U(O)\)-types indexed by \(\xi,\zeta\) each of length \(\le K\), sizes summing to \(k\), paired with \([\xi]\) and (after twisting) \([\zeta']\). The unitaries commute with inducing over locations, so an \(S_k\)-type \(\theta\) is compatible with these types only if \(c_{\xi,\zeta'}^\theta>0\), an LR coefficient (prime denotes conjugation). In this case
\begin{equation}
 0\le \theta_i-\xi_i\le K,\qquad
 0\le \theta'_j-\zeta_j\le K.                                         \label{ten:simultaneous:eq:4}
\end{equation}
In fact \(\xi\subseteq\theta\) and in an LR tableau of \(\theta/\xi\) each row has length at most the first count \((\zeta')_1\le K\). To see the row bound, let \(C_j\) be the counts before reading the row in the lattice word (read right to left), and \(d_j\) the counts in the row. Weak row increase and the lattice condition give \(d_j\le C_{j-1}-C_j\) for \(j\ge2\), so the length is at most \(d_1+C_1\). Conjugating the coefficient (sign twist) and using symmetry gives the column bound.

For fitting \(\theta\) set \(a_i=\theta_i,\ b_j=(\theta'_j-K)_+\) for \(1\le i,j\le K\); together these counts cover the cells. The hook formula gives
\[
 B_\theta\le H(a,b)\le B_\theta+C K^2\log(2n).
\]
Indeed the hook product is at least \(\prod a_i!\prod b_j!\) by counting to the right in the top rows and downwards below them. For an upper bound, corner hooks (in the \(K\times K\) square) are at most \(n\). Arm hooks to its right have product in row \(i\) at most \(a_i!\) (add at most \(K\) to the right-count plus one there), and leg hooks below have product in column \(j\) at most \((b_j+K)!/K!\le b_j!(2n)^K\). Finally log multinomial for these counts is at most their \(H\) and within \(O(K\log(2n))\) of it by \(\log z!=z\log z-z+O(\log(2n))\) for \(0\le z\le n\). By \eqref{ten:simultaneous:eq:4}, each of \(a_i,b_j\) differs from the respective \(\xi_i,\zeta_j\) by at most \(K\). So \(H(\xi,\zeta)\) (concatenating the counts) is within \(C K^2\log(2n)\) of \(H(a,b)\) (per unit change of an integer \(z\), \(z\log z\) changes by at most \(1+\log(2n)\)). Weights of these unitary types are majorized by \(\xi,\zeta\) respectively: by the semistandard tableaux formula any \(i\) symbols fill at most \(\sum_{j\le i}\xi_j\) cells since columns are strict, and similarly for \(\zeta\). By convexity (majorization inequality) their combined \(H\) is at least \(H(\xi,\zeta)\). This proves the weight assertion, unchanged on rotating by the parity unitaries.

For \eqref{ten:simultaneous:eq:3}, for a compatible unitary type and \(k>0\) take the diagonal density matrix with entries \((\xi_i/k,\zeta_j/k)\) and conjugate uniformly by the parity unitaries. Before conjugation its tensor power acts by scalars on weights, with eigenvalue \(\exp(-H(\xi,\zeta))\) on the highest weight of that type (zero entries with exponent zero contribute 1). It preserves every unitary submodule (it is a linear combination of weight projections, which preserve them by torus averaging); the conjugates do also. Thus, in an orthogonal unitary-module decomposition, on each irreducible copy of the given type its conjugation average is scalar at least this eigenvalue divided by the type dimension (by the trace there and Schur's lemma), and on the other modules positive semidefinite. That dimension is at most \((2n)^{C K^2}\) by the Weyl dimension formula, and there are at most \((n+1)^{2K}\) types. The isotypic projection \(\Pi_\theta\) is bounded by the sum of the compatible unitary-type projections. Averaging the densities over these choices as well (absorbing their count in the coefficient), and using the upper estimate on \(H(\xi,\zeta)\), proves \eqref{ten:simultaneous:eq:3}. In size zero the assertions use the scalar tensor power and a parity-preserving density can be chosen arbitrarily.

\subsubsection{Overlap on a grid with missing cells}
Here is the overlap estimate on this tensor space that will apply within the recursion. Use an \(s\)-by-\(r\) grid, \(n=rs\ge2,\ K\le n\), with \(t\) arbitrary holes and \(M=n-t\) remaining slots, using \(H_M\). Let \(\mu\vdash M\) satisfy the \(K\)-hook condition. Index rows and columns together by \(g\). In each take a positive semidefinite group-algebra operator on its remaining slots supported on a single partition type \(\delta_g\), with trace on \([\delta_g]\) at most \(c_g\); we can assume the type occurs in the graded tensor power there. Let \(A\) be the product for rows, \(B\) the product for columns (using the subgroup actions in \(S_M\)). Put \(S=\sum_g B_{\delta_g}\). Then
\begin{equation}
 \operatorname{tr}_{H_M}(\Pi_\mu A\Pi_\mu B)
 \ \le\ w_\mu \Big(\prod_g c_g\Big)
 \exp\!\left(C((s+r+1)K^2\log(2n)+t)+\min(0,S-B_\mu)\right).          \label{ten:simultaneous:eq:5}
\end{equation}
One bound uses \(\|B\|\,\operatorname{tr}_{H_M} A\): the column norms are at most \(c_g\), and the row traces contribute at most \(c_g w_{\delta_g}\) each. This gives the bound without the \(\min\) term.

For another bound replace both operators by their isotypic projections times the norm bounds, using positivity with \(\Pi_\mu\) inserted (each replacement traces against a positive semidefinite operator by cyclicity). Write the products of projections as \(P_A,P_B\); after the factors \(\prod c_g\) the trace left to bound is \(\operatorname{tr}(P_A\Pi_\mu P_B\Pi_\mu)\). Replace first \(P_A\) by \eqref{ten:simultaneous:eq:3} in the rows (tensoring the bounds), giving, apart from coefficients, an average of product states \(R\) with one density \(\rho_i\) used throughout the remaining cells of row \(i\). For each such term diagonalize \(\bar\rho=s^{-1}\sum_i\rho_i\) in a parity basis (extend by arbitrary parity-preserving densities for empty rows). We have \(\Pi_\mu\le Z=\sum_u Z_u\), the sum of weight projections in that basis obeying the weight assertion for \(\mu\). The positive \(P_B\) commutes with both \(\Pi_\mu\) and \(Z\) (centrality for the former, torus commutation for the latter), so \(\Pi_\mu P_B\Pi_\mu\le Z P_B Z\). Now similarly replace \(P_B\) giving product states \(Q\) with parity-preserving densities \(\sigma_j\) by column (extended for empty ones). Signed reorderings back to the joint slot order give the ordinary product states as the slot matrices preserve parity. The row and column coefficients from \eqref{ten:simultaneous:eq:3} together cost \(\exp(S+C(s+r)K^2\log(2n))\).

We bound \(\operatorname{tr}(R Z Q Z)=\|R^{1/2}Z Q^{1/2}\|_{\rm HS}^2\) using each \(Z_u\). This projection is the torus Fourier integral of \(g^{\otimes M}\) with unit-modulus multiplier, over diagonal phases \(g\). Put \(X=(\bar\rho+\varepsilon I)^{-1/2}\), \(\varepsilon>0\), with diagonal entries \(x_l\). In this integral we can insert \(X^{\otimes M}\) dividing by \(\prod_l x_l^{u_l}\), since \(X^{\otimes M} Z_u=(\prod_l x_l^{u_l}) Z_u\). By the norm triangle inequality we then use the integrand estimate
\[
 \|R^{1/2}(Xg)^{\otimes M} Q^{1/2}\|_{\rm HS}^2
   =\prod_{(i,j)\ \mathrm{remaining}}\operatorname{tr}(X\rho_i X g\sigma_j g^*)\le e^t .
\]
Indeed these nonnegative factors have sum over the entire extended grid at most \(n\), by averaging to \(X\bar\rho X\le I\) and the averaged trace-one column density. On the remaining grid apply arithmetic-geometric means and \(M\log(n/M)\le t\) (empty product if \(M=0\)). Also
\[
 \prod_l x_l^{-2u_l}\le (1+2K\varepsilon)^M \exp(-H(u)),
\]
by normalizing the entries of \(\bar\rho+\varepsilon I\) to probabilities (the product is maximized then at frequencies \(u_l/M\); for \(M=0\) the inequality is trivial). Let \(\varepsilon\) decrease to zero. There are at most \((n+1)^{2K}\) projections in the sum, so summing norms and squaring gives the penalty \(\exp(-B_\mu+C K^2\log(2n)+t)\). This proves the second bound. Together they imply \eqref{ten:simultaneous:eq:5}, where a factor \(w_\mu\ge1\) has been allowed.

\subsubsection{Recovering the marked trace recurrence}

We use the balanced grid and the arbitrary parent hook $\mu$ from
\eqref{ten:simultaneous:eq:6}. Thus $X$ is the product of row sweeps
and $Y$ the product of column sweeps, with $B_n=YX$; $g$ ranges over
both families of lines. The following calculation proves that recurrence
using the overlap just established.

First, by a trace power inequality in \([\lambda]\),
\[
 \mathcal L_\lambda\le
 \operatorname{tr}_\lambda\big((XX^*)^p (Y^*Y)^p\big).
\]
We use \(\operatorname{tr}((B^{1/2} A B^{1/2})^p)\le\operatorname{tr}(B^{p/2} A^p B^{p/2})\) with \(B=XX^*, A=Y^*Y\) and the polar decomposition for \(X\). For the powers of 2 this instance of the trace inequality can be seen as follows. For positive definite \(A,B\), products of the top \(j\) eigenvalues of \(AB\), squared, are bounded by products of the top \(j\) eigenvalues of \(A^2 B^2\). Indeed take the \(j\)-th exterior power of \(AB\), bounding spectral radius by operator norm, whose square gives the top eigenvalue of the exterior power of \(B A^2 B\). All eigenvalues used are positive and \(B A^2 B\) has the eigenvalues of \(A^2 B^2\). Iterate squaring; equality of the full products holds by determinant. Thus logs of the eigenvalues of \(AB\), multiplied by \(p\), are majorized by the logs for \(A^pB^p\); sum the exponentials. This proves the inequality and a limit gives the semidefinite case.

The powered elements on the right factor by groups because disjoint group operations commute within each part. Expand by multiplying in each group by its central isotypic projectors and summing. Each term has positive semidefinite row and column operators separately, products of group-algebra operators supported on \(\alpha_g\) with trace \(\mathcal L_{\alpha_g}\) on that type (either orientation of the power gives the same trace).

To bound a term, augment the even tensor alphabet for the \(K_n\)-construction by \(t\) distinct marks, each required to occur once, the rest of the slots from the usual \(V\). In addition keep only core type \(\mu\): in each placement \(z\) of the ordered marks use the isotypic subspace \(\Pi_\mu H_M\) in the unmarked slots, ordered by position. This direct sum \(\mathcal K\) over placements is invariant under signed \(S_n\); the even marks add no signs and each map of placements acts by the induced (signed) reordering on core slots. It is the induced module from the one-placement space for the mark stabilizer \(S_M\) (its translates over the cosets are exactly the placement summands). In particular the multiplicity of \(\lambda\) is \(w_\mu f^{\lambda/\mu}>0\) by reciprocity and branching. The term's trace in \([\lambda]\) is at most its trace on \(\mathcal K\) divided by that multiplicity, since all irreducible traces of a product of two positive semidefinite operators are nonnegative here.

In this last trace insert the placement-space projections between operators. Only diagonal placement blocks of each operator contribute: initial and intermediate placements in the two maps making a return must agree on each mark's row by the row operator, and on each mark's column by the column operator, hence agree entirely. On the diagonal block for placement \(z\), with \(h_g\) marks in group \(g\), we retain in each group-algebra factor just the coefficients on the subgroup fixing those slots pointwise. It acts on the \(m_g-h_g\) core slots there, commuting (as an operator in the whole core permutation algebra) with \(\Pi_\mu\). Thus this placement trace has \(\Pi_\mu\) inserted with both products of truncated operators on \(H_M\). Coefficient truncation to this subgroup preserves positivity, as seen by compressing the regular representation to the subgroup. It is supported only on types \(\delta_g\subseteq\alpha_g\): for an absent restriction type the original element times the subgroup type projector vanishes, and coefficient truncation commutes with right multiplication by subgroup elements. Its trace in \([\delta_g]\) is at most
\[
 c_g=\frac{D_{\alpha_g}\mathcal L_{\alpha_g}}{(m_g)_{h_g}D_{\delta_g}}.
\]
In fact regular trace is group size times the identity coefficient, so the subgroup regular trace equals the original \(D_{\alpha_g}\mathcal L_{\alpha_g}\) divided by \((m_g)_{h_g}\). Now sum over core types by using their central projectors in these subgroup elements, needing only types occurring on the tensor powers.

We may apply \eqref{ten:simultaneous:eq:5} on each placement for these terms. Substituting \(c_g\), the bound there has logarithm at most
\[
 \log w_\mu + \sum_g e_{\alpha_g} - \sum_g\log (m_g)_{h_g}
 -\max\big(B_\mu,\sum_g B_{\delta_g}\big)
 + C((s+r+1)K_n^2\log(2n)+t).
\]
For fixed counts, the placements number at most the product of the multinomial counts assigning rows and assigning columns to the \(t\) distinct marks. In each branch the multinomial is at most \(\exp(t\log t-\sum_{g\ \mathrm{in\ branch}} h_g\log h_g)\) (use category probabilities \(h_g/t\) if \(t>0\)). The log of the product together with the negative falling-factorial logs costs at most
\[
 2t\log t-\sum_g h_g\log h_g - t\log n+2t
 = \sum_g J(h_g,m_g)-2J(t,n)+2t .
\]
Here \((m)_h\ge (m/e)^h\) by the geometric mean of the largest \(h\) factors of \(m!\), and zero-log terms use the previous conventions. Counts of all vectors of partitions \(\alpha_g,\delta_g\) and counts vectors themselves have log at most \(C n^{3/4}\log(2n)\): a partition in size at most \(m_g\) can be specified by the lengths along the sides of its Durfee square (top rows and first columns) using at most \(2\sqrt{m_g}\) lengths plus the square size, and one more entry suffices for \(h_g\); \(r,s\) are within a constant factor of \(\sqrt n\). All the sublinear powers times logs in these errors are bounded by \(C n^\eta\).

Finally
\[
 D_\lambda\le \binom nt D_\mu f^{\lambda/\mu},
\]
by assigning entries of a standard tableau to the two regions and ignoring their boundary conditions. Also \(\log\binom nt\le J(t,n)+t\). Thus after division by \(w_\mu f^{\lambda/\mu}\), adding \(B_\lambda+J(t,n)\) to the log trace bound cancels the \(-2J(t,n)\) at cost at most \(B_\mu+t\). This gives \eqref{ten:simultaneous:eq:6}. Zero terms can be omitted; the choices needed in the max are nonempty when the parent trace is positive, by the trace bounds leading to it.
